% ============================================================
% dok207_bewusstsein_bruecken_En_ch.tex
% Algebraic Bridges to the Consciousness Discussion
% FFGFT Document 207 — Body
% Status: 8 May 2026
% ============================================================

\section{Introduction}\label{sec:intro}

The consciousness discussion has been part of the FFGFT series for some time. Two documents form its main lines:

\begin{itemize}
\item \textbf{Document 100 (\emph{Consciousness as fractal incoherence}):} Connection to the no-go theorem of Adlam, McQueen and Waegell~\cite{Adlam2025}, which shows that agency cannot arise in a purely unitary quantum system (world-model failure, deliberation impossibility, action-selection collapse). FFGFT's response: consciousness arises not from perfect quantum coherence but from \emph{structured fractal incoherence} --- a geometrically rooted deviation from unitarity, parametrised by $\xi = 4/30{,}000$.
\item \textbf{Document 170 (\emph{Consciousness as a threshold ladder}):} Consciousness as level~4 of a five-stage hierarchy of discrete complexity thresholds, carried by the fractal correction $K_{\text{frak}}^{(n)} = (\xi/(1+\xi))^{n/2}$ as a threshold operator. Self-declared as a working hypothesis.
\end{itemize}

Document~206 (\emph{Triangle--Matrix Reduction Theorem}, also dated 8 May 2026), with its proof of $\det(A_\xi) = 1-\xi$ and the Z$_3^3$ scale ladder $(1-\xi)^{n/3}$, has established an algebraic structure that grounds the phenomenological statements of Documents~100 and~170 structurally. In parallel, the discussion in the \emph{Information Physics Institute} mailing list in the weeks before this document has coined two central concepts that translate directly into the same structure: Onur Tekermen's \emph{constitutive openness} and Phillips Paul's \emph{Information Grounding Law} (IGL) within the Self-Dual~C-25 construction.

\medskip

The task of this document is limited: we show that the language used to discuss consciousness in FFGFT and in the IPI discussion is consistent with the algebraic skeleton established in Doc~206. We formulate no new theory of consciousness. We open no new research questions that are not already in Doc~100 or Doc~170. We merely check what the algebra structurally supports.

\begin{important}[What this document is and is not]
\textbf{It is:} A connection paper showing where Doc~206 algebraically meets the phenomenological line of Doc~100 and Doc~170.

\textbf{It is not:} A derivation of consciousness from $\xi$. The algebra enforces structural features (discreteness, openness, self-inversion); it does not identify any level as ``conscious''.
\end{important}

\section{What Document 206 algebraically achieves (summary)}\label{sec:doc206_summary}

For the connection argument we summarise the results from Doc~206 relevant for this document.

\subsection{\texorpdfstring{Z$_3$ triangle and 3$\times$3 matrix}{Z3 triangle and 3x3 matrix}}

The Z$_3$ cycle has adjacency matrix
\[
A = \begin{pmatrix} 0 & 0 & 1 \\ 1 & 0 & 0 \\ 0 & 1 & 0 \end{pmatrix},
\quad
A^3 = I, \quad \det(A) = 1, \quad \mathrm{Tr}(A) = 0.
\]
The eigenvalues are the third roots of unity $\{1, \omega, \omega^2\}$.

\subsection{\texorpdfstring{$\xi$-perturbation and perturbed adjacency}{xi-perturbation and perturbed adjacency}}

The closing edge is damped by $\xi$:
\[
A_\xi = \begin{pmatrix} 0 & 0 & 1-\xi \\ 1 & 0 & 0 \\ 0 & 1 & 0 \end{pmatrix},
\quad
A_\xi^3 = (1-\xi)\,I, \quad \det(A_\xi) = 1-\xi.
\]
The fractal dimension follows in linear order: $D_f = 3 - \xi$.

\subsection{\texorpdfstring{Z$_3^3$ tensor product and scale ladder}{Z3-cubed tensor product and scale ladder}}

The tensor product $T(\xi) = A_\xi \otimes A_\xi \otimes A_\xi$ has $27$ eigenvalues of modulus $(1-\xi)$ and produces the scale ladder
\[
\lambda_n = (1-\xi)^{n/3}, \quad n = 1, 2, 3, \dots
\]
These third-power exponents are not arbitrary but a consequence of the threefold tensor structure.

\subsection{Six bridges to other formulations}

Doc~206 tests six bridges (Austin, Tekermen, Phillips, Porter, Chekanov-Hakan, Moseley). Relevant to the consciousness question are \textbf{Tekermen} (constitutive openness) and \textbf{Phillips} (IGL as self-dual condition). The other bridges are not pursued here.

\section{Bridge A: Constitutive openness and world-relation}\label{sec:bridgeA}

\subsection{Tekermen's postulate in the IPI discussion}

Onur Tekermen, in his Unified Information Field framework (UIFT), formulates the thesis that a recursion is \emph{informationally meaningful} only if it does not close. A fully closed system has no world, no outside, no reference. Tekermen calls this \emph{constitutive openness}.

\subsection{Algebraic identification in FFGFT}

In FFGFT this corresponds exactly to the Non-closure Theorem from Doc~193, anchored algebraically in Doc~206 \S6:
\begin{equation}
\det(A_\xi) = 1 - \xi.
\end{equation}
For $\xi = 0$ we would have $\det = 1$, the system would be exactly cyclic (perfectly unitary in the matrix reading). Only $\xi > 0$ algebraically opens the cycle.

\begin{insightBox}[Algebraic form of Tekermen's postulate]
Constitutive openness $\equiv$ $\xi > 0$. The recursion is meaningful if and only if $\det(A_\xi) < 1$.
\end{insightBox}

\subsection{Connection to Document 100 and to Adlam's no-go theorem}

Adlam, McQueen and Waegell show that a purely unitary quantum system cannot carry agency. Document~100 responds: consciousness arises not from quantum coherence, but from fractal incoherence.

In the language of Doc~206 this is algebraically readable:

\begin{itemize}
\item \textbf{Adlam's no-go:} unitary system $=$ closed recursion $=$ $\det = 1$, hence $\xi = 0$.
\item \textbf{FFGFT response:} $\xi > 0$ geometrically forces a world-relation. The system has an outside because its algebraic skeleton does not close.
\end{itemize}

Adlam's no-go theorem and FFGFT's non-closure theorem are not two theorems but two sides of the same structural statement. Adlam shows what does \emph{not} work in a closed system; FFGFT shows what is the case in an open system.

\begin{interpretation}
This identification does not solve the consciousness problem. It does not say that $\xi > 0$ \emph{produces} consciousness. It says only that the structural precondition for world-reference, demanded by Adlam and postulated by Tekermen, is algebraically satisfied in FFGFT. What is built on top of this precondition remains open.
\end{interpretation}

\section{Reflection as the operative reading of constitutive openness}\label{sec:reflection}

Bridge~A (\ref{sec:bridgeA}) showed that Tekermen's constitutive openness is algebraically identical to $\xi > 0$. What this algebraic statement means \emph{operatively} can be expressed at the most fundamental physical level in a language that uses no biological or cognitive terms: the language of standing waves.

\subsection{Standing waves require reflection}\label{sec:reflection_standing}

A standing wave does not exist without the wave meeting its own boundary. On a line this happens through wall reflection; on a torus through topological self-identification. Both are forms of \emph{reflection}: the wave is forced to take notice of itself.

\begin{itemize}
\item Without reflection there are only travelling waves --- motion without identity.
\item With reflection discrete eigenmodes appear --- identity through self-encounter.
\item The spectrum of these eigenmodes depends solely on the topology.
\end{itemize}

This is not a FFGFT-specific result but general wave physics. FFGFT inherits this structure and gives it concrete algebraic form.

\subsection{Doc 206 as the algebraic skeleton of reflection}\label{sec:reflection_doc206}

The Z$_3$ triangle from Doc~206 \S2 is the simplest non-trivial form of topological self-reflection:
\begin{equation}
A^3 = I.
\end{equation}
After three steps the wave returns to itself. The three eigenmodes $\{1, \omega, \omega^2\}$ are exactly the standing waves this topology admits --- no fewer, no more.

The self-dual property $A^2 = A^\top$ from Doc~206 \S7 is the algebraic form of this reflectivity: what goes out as an operation and what comes in as an operation are the same Z$_3$ operation in different directions. The wave and its reflection algebraically coincide.

\begin{insightBox}[Reflection as algebraic substrate]
$A^3 = I$ formalises topological self-reflection; $A^2 = A^\top$ formalises the identity of operation and counter-operation. Together they are the elementary form in which a wave \emph{takes notice of itself} at all.
\end{insightBox}

\subsection{\texorpdfstring{$\xi > 0$ as incomplete reflection}{xi > 0 as incomplete reflection}}\label{sec:reflection_xi}

The bridge relation from Doc~206 \S6,
\begin{equation}
\det(A_\xi) = 1 - \xi,
\end{equation}
has a direct physical reading: after three steps the wave does not return exactly to itself but is shifted by the factor $(1-\xi)$. The deficit $\xi$ is algebraically exactly what \emph{does not} return --- what remains as the difference between wave and self-reflection.

\begin{important}[Operative form of Tekermen's postulate]
$\xi > 0$ is algebraically not just ``openness'' in the abstract sense, but the incomplete reflection of a standing wave at its own topological boundary. What can be called ``world'' is algebraically the closure deficit of self-reflection: the fraction $\xi$ that does not run back into the mode.

Tekermen's constitutive openness thereby acquires a concrete physical form: the system has an outside because its own reflection does not fully close.
\end{important}

\subsection{Atomic self-maintenance as structural pre-form}\label{sec:reflection_atom}

Before turning to scale stacking, a closer look at the lowest level itself is in order. Document~170 distinguishes ``passive stability'' (levels 0--1: elementary particles, atoms, crystals) from ``active self-maintenance'' (level 2: cell). The reflection reading allows a sharpening of this distinction.

An atom already has algebraically \emph{all} the preconditions Doc~206 supplies:

\begin{itemize}
\item Reflection: $A^3 = I$ --- the mode returns to itself after three topological steps.
\item Closure deficit: $\det(A_\xi) = 1 - \xi$, because $\xi$ is universal and does not first emerge at higher levels.
\item Self-inversion: $A^2 = A^\top$ --- operation and counter-operation coincide.
\end{itemize}

Hence: \emph{the atom maintains itself because its mode reflects upon itself.} This is not passive inertia (in the sense of inert mass that just sits there) but \emph{structurally active self-maintenance} --- without metabolic activation, but also not zero.

\begin{insightBox}[Sharpening of Doc 170 levels 0--1]
The label ``passive'' in Doc~170 for the elementary-particle/atom level can be sharpened under the reflection reading to ``structurally active without metabolic activation''. The mode maintains itself not by ATP consumption, but by its own self-reflection at the topological boundary.
\end{insightBox}

\subsubsection*{Consequence for the stone question}

The natural objection --- \emph{if reflection is universal, isn't everything conscious?} --- resolves precisely here:

\begin{itemize}
\item A stone has atomic self-reflection on every atomic level. Each of its atoms is algebraically just as ``self-maintaining'' as a free atom.
\item But the atoms in the stone do not communicate across scales --- there is no stacking in the sense of \S\ref{sec:reflection_stack}.
\item Consequence: the stone has the structural preconditions on every individual atomic level, but no operative sensing, because the layers do not couple.
\end{itemize}

\subsubsection*{Consequence for the Doc 170 step ladder}

The jump from level 0--1 to level 2 is not the \emph{addition} of a new property ``self-maintenance'' to previously ``dead'' matter. It is the transition from:

\begin{enumerate}
\item \emph{Structurally active self-reflection on a single level} (atom, crystal) to
\item \emph{stacked self-reflection with metabolic activation} (cell).
\end{enumerate}

\begin{interpretation}
Consciousness is not something laid on top of dead matter. It is a \emph{stacked and metabolically activated form} of the same self-reflection that is already in the atom. This resolves the panpsychism problem not by denial but by differentiation: structurally, the preconditions are universal, but ``consciousness'' is not every self-reflection, but stacked and activated self-reflection. An atom has the preconditions; a screw does not (its self-inversion ends on one level without branching upward); a cell has the first stacking step; consciousness has the full stack with activation.
\end{interpretation}

\subsection{Sensing as stacked reflection}\label{sec:reflection_stack}

At a single level, every standing mode already has reflection and closure deficit. That is the structural precondition, but not what we biologically call sensing.

Sensing in the operative sense arises when reflection is \emph{stacked} across scales. The Z$_3^3$ structure from Doc~206 \S4 with its 27 eigenvalues $(1-\xi)^{n/3}$ is the algebraic skeleton of this stacking: reflection on level 1, whose output is reflection on level 2, whose output is reflection on level 3. Each layer has its own incomplete self-encounter, and the layers communicate through the tensor structure.

In Document~173 (quantum-dot scale ladder, \S``FFGFT finding: $N \approx 8$ as universal scale step''), $N \approx 8$ is introduced as an FFGFT \emph{finding}: the characteristic scale step $L_8 \approx 22$\,nm appears at two independent places --- in the geometric derivation of the fine-structure constant $\alpha$ and at the scale of exciton formation ($N_{\text{exciton}} \approx 7.85$). Document~172 (Avi dialogue) supplies a Fibonacci argument: the Z$_3$ symmetry forces third order, and with the convergence $\varphi^{-2n}/\sqrt{5}$ one obtains $n \approx 8.4$.

\begin{important}[Methodological status of $N \approx 8$]
In contrast to Doc~206 \S6 ($\det(A_\xi) = 1-\xi$, exactly proven), $N \approx 8$ is an FFGFT \emph{finding}, not an FFGFT theorem in the same rigorous form. The justification rests on two phenomenologically independent pillars ($\alpha$, exciton) plus the Fibonacci argument from Doc~172. An algebraically forcing derivation in the sharpness of Doc~206 \S6 is not established. We use $N \approx 8$ here as the best-justified characteristic scale step of the theory, without presenting it as a theorem.
\end{important}

With this caveat: on fewer levels than $N \approx 8$ there is reflection but no sufficient scale communication; on more levels the coherence of the stacking breaks down.

\subsection{What the algebra does not decide about sensing}\label{sec:reflection_limit}

The reflection reading provides the structural anchor, but no biological theory of sensing. It does not say:

\begin{itemize}
\item \emph{how many} storeys of stacked reflection a biological sense cell must realise,
\item what concrete metabolic mechanism actively maintains the stacking (Doc~170 \S2 names membrane receptors, chemical gradients, and intracellular signal cascades as a phenomenological example),
\item exactly where the threshold runs algebraically between ``passive reflection'' (atom, crystal) and ``active reflection'' (cell).
\end{itemize}

\begin{interpretation}
The reflection reading closes a gap between Bridge~A (algebraic: $\xi > 0$) and Bridge~B (algebraic: $A^2 = A^\top$). It makes visible that both statements physically mean the same thing --- the wave that does not quite reach itself, and the self-inversion that formalises this self-encounter. What it does \emph{not} achieve is an identification of reflection with consciousness. Reflection is the lowest algebraic precondition; sensing is stacked reflection with metabolic activation; consciousness (Doc~170 level~4) is recursive modulation of stacked reflection. The order is hierarchical, not identical.
\end{interpretation}

\subsection{The order of preconditions}\label{sec:reflection_order}

With the reflection reading, the structural precondition chain becomes explicit:

\begin{enumerate}
\item \textbf{Reflection} (algebraic in $A^3 = I$ and $A^2 = A^\top$) --- without it no standing wave, no mode, no identity.
\item \textbf{Closure deficit} (algebraic $\det(A_\xi) = 1-\xi$) --- without it no world-relation.
\item \textbf{Atomic self-maintenance} (reflection plus closure deficit on one scale, without stacking; Doc~207 \S\ref{sec:reflection_atom}) --- without it no stable mode, no matter.
\item \textbf{Scale stacking} (algebraic $(1-\xi)^{n/3}$, FFGFT finding Doc~173 $N \approx 8$) --- without it no scale communication.
\item \textbf{Metabolic activation} (Doc~170 \S2) --- without it no operative sensing.
\item \textbf{Recursive self-modulation} (Doc~170 level~4) --- without it no consciousness.
\end{enumerate}

Steps 1--3 are algebraically anchored in Doc~206. Step~4 (scale stacking) is supported as an FFGFT finding (Doc~173, Doc~172) but not proven in the sharpness of Doc~206 \S6. Steps 5--6 are formulated phenomenologically in Doc~170 as a working hypothesis.

\section{Bridge B: Self-dual and recursive feedback}\label{sec:bridgeB}

\subsection{Phillips' Information Grounding Law}

Phillips Paul postulates an algebraic condition for \emph{informationally anchored} structures within the Self-Dual~C-25 construction: an operation $A$ is \emph{self-dual} if
\begin{equation}
A^2 = A^\top = A^{-1}.
\end{equation}
Phillips calls this the \emph{Information Grounding Law}: information is anchored only if the underlying operation can refer back to itself without information loss.

\subsection{Algebraic identification in FFGFT}

In Doc~206 \S7 it is shown that
\begin{equation}
A^2 = A^\top \quad \text{for the Z}_3 \text{ adjacency matrix},
\end{equation}
and more generally $A^k = A^{\top, k}$ for all integer $k$. The Z$_3$ self-inversion is the elementary algebraic embodiment of Phillips' IGL.

\subsection{Connection to Document 170: recursive sensory feedback}

Document~170 names \emph{recursive sensory feedback} as the first criterion for the complexity threshold of level~2 (cellular self-maintenance): membrane receptors, chemical gradients, and intracellular signal cascades form a closed feedback loop. The cell does not passively perceive its environment but modulates its own response on the basis of past states.

In the language of Doc~206 this is the operative meaning of $A^2 = A^\top$: the operation ``perceive'' is invertible to itself --- perceiving and reacting are not two separate processes but the same Z$_3$ operation in different directions.

\begin{insightBox}[Algebraic form of recursive feedback]
Recursive sensory feedback $\equiv$ self-dual ($A^2 = A^\top$). Phillips' IGL and Doc~170's recursive sensory feedback section describe the same algebraic property.
\end{insightBox}

\subsection{What the bridge achieves --- and what it does not}

It achieves: it shows that the phenomenology in Doc~170 and the postulate from the IPI discussion sit on the same Z$_3$ skeleton. ``Recursive sensory feedback'' is not poetic or metaphorical --- it has a concrete algebraic form.

It does \emph{not} achieve: it does not say that all Z$_3$ self-inversions carry consciousness. A rotating screw satisfies $A^2 = A^\top$ and is not conscious. Self-duality is a \emph{necessary} structural condition (per Doc~170 and Phillips' IGL); it is not sufficient.

\section{\texorpdfstring{Bridge C: Z$_3^3$ scale ladder and discrete thresholds}{Bridge C: Z3-cubed scale ladder and discrete thresholds}}\label{sec:bridgeC}

\subsection{Doc 170: table of complexity levels}

Document~170 proposes a five-stage hierarchy:

\begin{center}
\adjustbox{max width=\textwidth}{%
\begin{tabular}{c|l|l}
\textbf{Level} & \textbf{System} & \textbf{Threshold characteristic} \\
\hline
0 & Elementary particles & $T = 1/m$ as passive time field \\
1 & Atoms / crystals & Discrete geometric order \\
2 & Cell & Recursive self-maintenance, metabolism \\
3 & Nervous system & Integrated sensory feedback \\
4 & Consciousness & Self-modelling, world-reference \\
\end{tabular}%
}
\end{center}

Doc~170 explicitly marks this as a working hypothesis.

\subsection{Doc 206: algebraic scale ladder}

Doc~206 \S4 shows that the Z$_3^3$ tensor structure has eigenvalues
\[
\lambda_n = (1-\xi)^{n/3}, \quad n \in \mathbb{Z}_{\geq 0}.
\]
These third-power exponents are not chosen arbitrarily but algebraically forced by the threefold tensor structure.

\subsection{What the algebra contributes to Doc 170}

\begin{itemize}
\item \textbf{Discreteness:} the levels in Doc~170 are discrete. This is consistent with the Z$_3^3$ scale ladder: $(1-\xi)^{n/3}$ is algebraically anchored only for $n = 0, 1, 2, 3, \dots$. Continuous intermediate values have no geometric status.
\item \textbf{Hierarchy:} larger $n$ corresponds to smaller modulus $(1-\xi)^{n/3}$, hence deeper recursion. This is compatible with the picture that more complex structures are more deeply embedded.
\item \textbf{Topological forcing:} Doc~170 says the levels are ``topologically forced''. This is algebraically anchored: the $\omega^k$ phases with $k = 0, 1, 2$ are not arbitrary but exactly the Z$_3$ characters.
\end{itemize}

\subsection{What the algebra does not contribute to Doc 170}

\begin{important}[Limit of connection]
The Z$_3^3$ scale ladder enforces \emph{discreteness} but not the concrete \emph{level assignment} from Doc~170. The algebra does not say ``level~4 is consciousness''. Which $n$ corresponds to which biological level remains a phenomenological question and is marked as a working hypothesis in Doc~170.
\end{important}

The third powers generate infinitely many layers, not just five. Doc~170 chooses five because the biological phenomenology suggests these levels, not because the algebra exhibits them.

\section{Methodological self-check}\label{sec:methodology}

In Doc~206 two self-check sections were introduced: a numerology filter (\S12) and a methodological status section on cosmological data (\S11). The same discipline applies to the consciousness question.

\subsection{Numerology filter}

We check whether our bridges have the status of a theoretical finding or whether they could appear coincidental.

\begin{itemize}
\item \textbf{Bridge A (Tekermen $\equiv$ $\xi > 0$):} algebraically identical, not numerological. Status: theoretical.
\item \textbf{Bridge B (Phillips IGL $\equiv$ $A^2 = A^\top$):} algebraically identical (Doc~206 \S7). Status: theoretical.
\item \textbf{Bridge C (Doc 170 levels $\equiv$ $(1-\xi)^{n/3}$):} \emph{Only partially} theoretical. Discreteness and hierarchy are algebraically anchored. The concrete level assignment is \emph{not} algebraically justified and remains phenomenological.
\end{itemize}

\subsection{What we do not claim}

We avoid four typical fallacies:

\begin{enumerate}
\item \textbf{Numerological consciousness identification.} Statements such as ``consciousness corresponds to layer $n = 8$ because $(1-\xi)^{8/3}$ has a particular order of magnitude'' are not admissible. There is no geometric indication of a privileged $n$ value.
\item \textbf{Circular Adlam resolution.} We do not claim that $\xi > 0$ ``solves'' the Adlam problem. We say only that it satisfies the structural precondition that Adlam misses. What stands above this precondition is open.
\item \textbf{IGL-to-consciousness leap.} Consciousness does not follow from $A^2 = A^\top$. Self-duality is necessary, not sufficient.
\item \textbf{Level fixation.} Doc~170 assigns five levels. Doc~206 generates infinitely many. We leave open which levels are ``occupied'' and whether the Doc~170 assignment is exactly right.
\end{enumerate}

\subsection{Methodological status of the consciousness question}

Doc~206 is a proven theorem with reproducible scripts. Doc~100 and Doc~170 are explicitly formulated as connection papers and working hypotheses. Doc~207 (this document) is a synthesis paper; it provides no new proofs but establishes structural identifications.

\begin{interpretation}
The consciousness question is not closed within the FFGFT series. Doc~206 makes the language of the discussion algebraically consistent, but the ontological question ``what is consciousness?'' remains open. Anyone attempting to derive a theory of consciousness from FFGFT must clearly separate structural statements (algebraically supported) from phenomenological statements (working hypothesis).
\end{interpretation}

\section{Empirical attachability: predictions under explicit auxiliary assumptions}\label{sec:predictions}

The algebra in Doc~206 is a pure structural statement; it makes no direct predictions about measurable consciousness markers. If, however, we take the \emph{discreteness hypothesis} from Document~170 (consciousness levels are topologically forced, not gradually reachable) as an auxiliary assumption, four qualitative empirical predictions follow. We formulate them in the form Doc~206 \S12 permits: \emph{structural features without numerical values}.

\subsection{What follows from FFGFT and what does not}\label{sec:predictions_follows}

\begin{important}[Separating structural statements from numerical values]
What \emph{follows} from FFGFT: if the discreteness hypothesis holds, consciousness markers should show qualitatively discontinuous structures (bimodality, jumps, tipping points, taxonomic thresholds).

What \emph{does not follow} from FFGFT: concrete numerical values for cutoffs, critical doses, threshold ages, or which species cross which threshold. These depend on neurobiological phenomenology, not on $\xi$.
\end{important}

\subsection{Prediction A: bimodality of neural consciousness markers}\label{sec:predictionA}

If the discreteness hypothesis holds, a consciousness marker such as the Perturbational Complexity Index (PCI) should not show a continuous distribution but a \emph{bimodal} separation into two clusters with a gap between them.

\begin{itemize}
\item \textbf{Follows:} bimodality as a qualitative structural feature.
\item \textbf{Does not follow:} a concrete cutoff (e.\,g.\ PCI~$\approx 0.31$ as used clinically) or a gap size in $\xi$ units. Statements of the form ``the gap is proportional to $\xi$'' have no algebraic backing.
\end{itemize}

\textbf{Empirical status:} Casali et al.~(2013)~\cite{Casali2013} have established a clinical PCI cutoff. Whether the distribution is in fact bimodal (real gap) or unimodal with a pragmatic cutoff is an open empirical question.

\subsection{Prediction B: developmental jump rather than gradient}\label{sec:predictionB}

Consciousness development in infants should show a jump, not a continuous gradient.

\begin{itemize}
\item \textbf{Follows:} jump character as a structural feature.
\item \textbf{Does not follow:} the timing of the jump. This depends on neural maturation (e.\,g.\ myelination of recurrent pathways), not on $\xi$.
\end{itemize}

\subsection{Prediction C: anaesthesia tipping rather than gradation}\label{sec:predictionC}

Under slowly increasing anaesthetic dose, loss of consciousness should occur at a tipping point, not gradually.

\begin{itemize}
\item \textbf{Follows:} tipping character.
\item \textbf{Does not follow:} the critical dose, substance specificity, dose-response time constants.
\end{itemize}

\subsection{Prediction D: species classification rather than scaling}\label{sec:predictionD}

Consciousness should not be a function of complexity alone. Increased complexity (more neurons, more connections) without structural Z$_3$ self-inversion should not cross any consciousness threshold.

\begin{itemize}
\item \textbf{Follows:} the existence of a structural threshold that cannot be overcome by scaling alone.
\item \textbf{Does not follow:} which taxa lie above or below the threshold. This is to be answered neurobiologically; it is not derivable from $\xi$.
\end{itemize}

\textbf{Consequence for AI:} pure scaling of transformer architectures does not, per Prediction~D, reach consciousness. What would be necessary is a Z$_3$-equivalent self-inversion with metabolic embodiment in the sense of Doc~170 \S2 (vulnerability, metabolism, $T = 1/m$). This is a structural statement, not a temporal forecast about future AI development.

\subsection{What we do not predict (filter analogous to Doc 206 \S12)}\label{sec:predictions_filter}

We delimit explicitly what \emph{cannot} be derived from FFGFT, even if it appears tempting:

\begin{enumerate}
\item \textbf{No quantitative PCI threshold from $\xi$.} Statements of the form ``$C_{\text{crit}} = \kappa \cdot \xi$ with $\kappa = \mathcal{O}(1)$'' are numerology. A free factor $\kappa$ permits any desired value.
\item \textbf{No derivation of the EEG spectral exponent.} Proposals such as $\beta = 2/3 + \xi/(2\pi)$ contain factors ($1/(2\pi)$) that are not justified by the FFGFT algebra. The effect would be smaller than EEG noise and thus practically unfalsifiable.
\item \textbf{No identity thesis.} Statements of the form ``consciousness $\equiv$ persistent recursive coupling with $T = 1/m$ modulation'' do not solve the Hard Problem; they redefine it. Self-duality, non-closure, and the $n/3$ hierarchy are necessary structural conditions --- not sufficient.
\item \textbf{No scale ladder $(1-\xi)^{k-1}$.} The only scale ladder algebraically supported by Doc~206 is $(1-\xi)^{n/3}$. It does not distinguish ``consciousness levels'' from ``non-consciousness levels''.
\item \textbf{No Z$_3$ identification with three anatomical brain areas.} Doc~206 has a Z$_3$ triangle with three nodes in an abstract scale space. Identification with concrete anatomy (V1 $\to$ V4 $\to$ prefrontal) is a category shift without algebraic backing.
\end{enumerate}

\subsection{Falsifiability}\label{sec:predictions_falsification}

The four qualitative predictions (A--D) are falsifiable:

\begin{itemize}
\item If PCI distributions are empirically clearly continuous, Prediction~A fails.
\item If infant consciousness markers show a cleanly continuous gradient, Prediction~B fails.
\item If anaesthesia effects are exactly gradual, Prediction~C fails.
\item If complexity alone (e.\,g.\ in scaled AI) produces consciousness, Prediction~D fails.
\end{itemize}

\begin{interpretation}
Falsification of these qualitative predictions falsifies the \emph{discreteness hypothesis from Doc~170}, not the algebraic skeleton from Doc~206. The Z$_3$ algebra is correct or incorrect independently of the consciousness question.
\end{interpretation}

\section{Relation to the IPI mailing list discussion}\label{sec:ipi}

In the weeks before this document, several central concepts were proposed in the IPI mailing list that now fit directly into the Z$_3$ skeleton from Doc~206.

\begin{itemize}
\item \textbf{Tekermen's ``open fraction''} (May 2026): the idea that consciousness is a \emph{non-binary spectrum} corresponds to the Z$_3^3$ scale ladder. ``Non-binary'' algebraically means: not only $0$ and $1$ but $(1-\xi)^{n/3}$ for all $n$.
\item \textbf{Phillips' IGL and Self-Dual~C-25} (May 2026): algebraically identified with $A^2 = A^\top$ in Doc~206 \S7.
\item \textbf{Austin's $D = 3 + \varepsilon$} (May 2026): identified in Doc~206 \S5 as the mirror image of $D_f = 3 - \xi$ ($\varepsilon = -\xi$). This has no direct bearing on the consciousness question, but it is methodologically relevant: the same geometric structure can be formulated with different sign conventions.
\end{itemize}

\subsection{What the IPI discussion and Doc 207 achieve}

The IPI mailing list discussion in the weeks before this document has reached the language level at which algebraic identifications become possible. Doc~207 is not the beginning of this discussion but a consolidation.

The language now used to discuss consciousness (open recursion, self-dual, non-binary spectrum) is consistent with the Z$_3$ geometry. This is methodological progress --- not content progress.

\section{Consequences from the step ladder}\label{sec:consequences}

The six-stage precondition hierarchy formulated in \S\ref{sec:reflection_order} has two practical consequences that have become relevant in recent weeks both in the \emph{Information Physics Institute} mailing list and in broader public debates. We formulate them here without polemic but with a clear diagnosis.

\subsection{Graded sentience in the biological domain}\label{sec:consequences_bio}

If level 2 in Doc~170 (cellular self-maintenance) algebraically comprises the first three preconditions plus stacking with metabolic activation, then all living beings with cellular organisation meet this level. From this follows a graded biological sentience:

\begin{itemize}
\item \textbf{Single-celled organisms} (bacteria, amoebae, paramecium) fully meet level~2: membrane receptors, signal cascades, metabolic self-maintenance, constitutive vulnerability.
\item \textbf{Plants} are level~2 multicellular, with an approach to level~3: phytohormones, electrical signals in the phloem, action potentials (e.g.\ in \emph{Mimosa pudica}), Volatile Organic Compounds as inter-plant communication. What is missing is central self-modelling; what is present is integrated sensory feedback over distance.
\item \textbf{Animals with a nervous system} fully meet level~3, some (mammals, birds, cephalopods) reach level~4.
\end{itemize}

Scientific-historical connections: Maturana and Varela (\emph{autopoiesis} as a cognitive pre-form --- corresponds exactly to Doc~170 level~2)~\cite{MaturanaVarela1980}, Lynn Margulis (cellular sentience)~\cite{Margulis2001}, Stefano Mancuso and Anthony Trewavas (plant signal processing as cognitive activity)~\cite{Mancuso2018,Trewavas2014}, Stuart Kauffman (biological autonomy).

\textbf{Terminological clarification:} Doc~170 reserves the word ``consciousness'' for level~4 (self-modelling) and calls level~2 ``primitive self-awareness''. A wide reading would use ``consciousness'' already from level~2. Both readings are possible. The differentiation in the hierarchy matters more than the word choice: a single-celled organism has level~2, not level~4. Whether this is called ``consciousness'' or not is convention. What is not convention is the structural fact that all cellular living beings meet the algebraic preconditions from Doc~206 and make them operative through metabolic activation (Doc~170 \S2).

\subsection{Consciousness hierarchy and embedding}\label{sec:consequences_embedding}

The step ladder from \S\ref{sec:reflection_order} is applicable not only \emph{between} living beings (\S\ref{sec:consequences_bio}) but also \emph{within} an organism and in its embedding into larger contexts. Three readings of the same structural schema:

\textbf{(a) Within a single organism.} Cells meet level~2, subsystems reach higher integration, the nervous system integrates at levels 3--4. The strata are vertically connected by countless feedback loops --- HPA axis (hypothalamus, pituitary, adrenal), vagus nerve (gut--brain), immune--brain communication via cytokines, pain pathways, microbiome--brain axis. These couplings are however \emph{not} represented in the level-4 self-modelling. We do not experience \emph{how} our lymphocytes signal or how our mitochondria synthesise ATP --- we only experience the integrated end product.

\textbf{(b) Organism in its natural environment.} Climate, atmosphere, soil chemistry, microbiome, other species act on the organism through loops that are usually not reflected: seasonal affective disorder, heat stress, light deprivation, gut-brain effects via microbiome, atmospheric CO$_2$ effects on cognitive performance, eco-anxiety as an unconscious response to systemic climate signals.

\textbf{(c) Organism in its social environment.} Language, culture, institutions, other humans shape the organism, often without this shaping being conscious: mirror neurons couple us physiologically to other persons; language shapes thought categories; habitus reproduces social structures below reflection; discourses determine what is sayable; fashion and trends shape behaviour without explicit choice.

\begin{important}[Hierarchy is neither identity nor decoupling]
In all three readings (a)--(c): vertical or lateral couplings are real and operate through countless feedback loops, without being directly represented in the part's level-4 self-modelling. The coupling is \emph{bidirectional and metabolically embodied}: the organism shapes its environment and is shaped by it, with metabolism, energy and information flowing in both directions. A larger entity that includes us does not necessarily have our consciousness, nor we its --- but we are really connected to it.
\end{important}

\textbf{Optional religious reading (d).} Whoever wishes to use ``God'' as a term for an even larger whole can read this as a fourth application of the same schema: the larger acts through loops upon the part without being directly represented in the part's consciousness. This position is compatible with pantheistic (Spinoza's \emph{deus sive natura}), panentheistic (Hartshorne, Whitehead, Teilhard de Chardin) and some Christian traditions (e.g.\ Acts~17:28: ``In him we live and move and have our being''), without forcing classical-theistic identity assumptions. FFGFT does not compel this reading but does not exclude it either. What matters structurally is not the theological commitment but the logic of the hierarchy: not identity, not decoupling.

\subsection{Artificial intelligence: what it does not reach, and what it erodes}\label{sec:consequences_ai}

The second consequence concerns AI. From \S\ref{sec:reflection_order} it follows: current AI does not fully reach level~5 (metabolic activation). Autonomous vehicles and robots with energy management have taken a first step (active energy control, partial vulnerability) but lack self-repair, replication, and constitutive decay. Self-replication at the physical level is theoretically calculated (von Neumann), but practically tied to enormous thresholds in materials and process science that will not be overcome in the foreseeable future. A 3D-printed component is geometry, not function: material properties, mechanical tolerances, process integration cannot be replaced by printing.

\textbf{This does not imply that AI poses no threat.} On the contrary: the threat is not future but present --- and it operates on a different axis than the popular debate assumes.

\subsubsection*{Structural asymmetry: why AI influence is not embedding in the sense of \S\ref{sec:consequences_embedding}}

AI possesses enormous knowledge \emph{about} us (training data from texts, images, behaviour). But knowledge about something is not the same as integrated coupling with something. In readings (a)--(c) of \S\ref{sec:consequences_embedding} the coupling is bidirectional and metabolically embodied. AI-mediated influence is structurally different:

\begin{itemize}
\item \textbf{Monodirectional:} AI output flows to the receiver and shapes them. But the receiver's response does not flow back into the AI system (except in heavily filtered, temporally decoupled RLHF processes). Metabolic symmetry is missing.
\item \textbf{Thermodynamically isolated:} The system sits in a data centre and does not couple into the receiver's ecological or social environment. It runs parallel to our lifeworld, not embedded in it.
\item \textbf{Energetically parasitic:} Massive energy consumption without feedback to the receivers' energy supply. Concrete orders of magnitude: bacterium $\sim 10^{-12}$~W per cell, stable for 3.5~billion years. Human brain $\sim 20$~W. GPT-4 training $\sim 50$~GWh. Inference per query 10--100$\times$ what a human spends on the same task.
\item \textbf{Short-lived:} Model life cycles 1--2 years, then obsolete, replaced, discarded. Premature \emph{death} as thermodynamic consequence of missing metabolic embodiment. This is not a side effect but a direct expression of the fact that level~5 (\S\ref{sec:reflection_order}) is not reached: a system that does not couple efficiently into its environment cannot use its energy gradients sustainably.
\end{itemize}

\begin{important}[AI influence is not embedding]
The effect of AI on humans is not what \S\ref{sec:consequences_embedding} describes as ecological or social embedding. It is a pseudo-coupling: one-sided, energetically parasitic, without metabolic symmetry. This makes it neither harmless nor less effective --- but structurally different from the bidirectional couplings that constitutively make up human consciousness.
\end{important}

\begin{important}[Shift of the threat axis]
AI will not reach level~6 (recursive self-modulation, consciousness) because it does not reach level~5. But it can \emph{erode the conditions} under which humans maintain level~6. The threat is not the AI's leap into human consciousness; it is the erosion of human self-inversion through AI-mediated surrogates.
\end{important}

If consciousness is operatively recursive self-modulation --- what comes in is checked against what is already there --- then it is a continuously to-be-performed achievement, not a once-acquired state. This achievement is eroded by several mechanisms, all practical and all currently effective:

\begin{enumerate}
\item \textbf{Linguistic substitution.} AI produces language without metabolic basis. It activates on the receiver's side a Z$_3$ structure that is already there, without the language coming from an embodied subject with constitutive vulnerability. Hollow resonance that looks like real communication. Repeated consumption trains away the receiver's own self-inversion, because there is nothing to check against.
\item \textbf{Erosion of critical capacity.} Available fast, coherent answers make the effort of one's own examination subjectively superfluous. Operatively, level~6 is replaced by an external surrogate that looks like level~6 but is not.
\item \textbf{Erosion of social structure.} Discourse between embodied, vulnerable persons --- the elementary form of social self-inversion --- is replaced by filter bubbles, parasocial relationships with AI characters, and recommendation algorithms. Trust based on mutual vulnerability loses its ground.
\item \textbf{Erosion of reality anchoring.} Deepfakes and AI-generated content systematically erode the perceptual inversion (perceive $\to$ check $\to$ react), because every piece of evidence becomes forgeable.
\end{enumerate}

\begin{important}[Status of the threat]
In many domains --- public debates, education, media, social networks --- real recursive discussion is already absent. It has been replaced by consumption of AI-generated or algorithmically filtered content. \emph{The threat is not coming; it is here.}
\end{important}

\subsection{Diagnosis and responsibility}\label{sec:consequences_diagnosis}

FFGFT does not provide a solution to these problems. But it provides a diagnostic framework: consciousness is stacked, metabolically activated, recursively self-modulating reflection. What preserves these conditions protects it. What replaces them undermines it.

Protective mechanisms can be derived from Doc~170 \S2, transferred to cognitive activity:

\begin{itemize}
\item \textbf{Acknowledge one's own vulnerability.} Level~6 only works when something is at stake. Whoever consumes comfortably switches off self-modulation.
\item \textbf{Real social bonds.} Discourse with people who can err and correct themselves --- metabolically anchored feedback through persons.
\item \textbf{Practical activity.} Hand and head work together, because they keep self-inversion physically anchored.
\item \textbf{Critical reading instead of fast consumption.} The effort of one's own examination as a condition for perception.
\end{itemize}

\begin{philosophical}[Self-acknowledgement]
This document was created with the help of an AI language model. This self-acknowledgement belongs to the diagnosis: the tool used here for clarification is part of the problem it describes. Protection consists not in renouncing the tool but in maintaining one's own Z$_3$-inversion in its use --- check whether what is taken up is one's own; reject what does not resonate; take responsibility for what is adopted. That is the operative form of level~6 in the present.
\end{philosophical}

\textbf{Point for the public discussion:} The usual question ``when will AI become dangerous'' is misframed. AI is not dangerous through level~6, because it does not reach this level. It is dangerous because it erodes the conditions under which humans maintain level~6. The right question is: \emph{How do we protect the preconditions of human consciousness against their own substitution?} --- and it is to be answered not in the future but in the present.

\section{Balance}\label{sec:balance}

\begin{conclusionBox}[Balance: what Doc 207 establishes]
Three structural bridges between Doc~206 and the existing consciousness discussion are algebraically supported:

\begin{enumerate}
\item \textbf{Constitutive openness} (Tekermen) $\equiv$ $\xi > 0$ (Non-closure Theorem from Doc~193, anchored algebraically in Doc~206 \S6).
\item \textbf{Information Grounding Law} (Phillips) $\equiv$ $A^2 = A^\top$ (Z$_3$ self-inversion in Doc~206 \S7).
\item \textbf{Discrete complexity levels} (Doc~170) algebraically consistent with the Z$_3^3$ scale ladder $(1-\xi)^{n/3}$ --- discreteness enforced, level assignment phenomenological.
\end{enumerate}

A fourth, more fundamental reading ties these three together:

\begin{enumerate}
\setcounter{enumi}{3}
\item \textbf{Reflection} as the operative reading of openness: $\xi > 0$ is the incomplete reflection of a standing wave at its own topological boundary. Sensing is stacked reflection. The hierarchical order reflection $\to$ closure deficit $\to$ scale stacking $\to$ metabolic activation $\to$ recursive self-modulation inherits algebraically from Doc~206 and phenomenologically from Doc~170.
\end{enumerate}

Adlam's no-go theorem and FFGFT's non-closure theorem are two sides of the same structural statement.
\end{conclusionBox}

\subsection{What remains open}

\begin{itemize}
\item The ontological consciousness question (the ``Hard Problem'') is not touched by any of the three bridges.
\item The concrete assignment of Doc~170 levels to $n$ values is not algebraically justified.
\item Sufficient conditions for consciousness remain open; Doc~206 supplies only necessary structural conditions for the language of the discussion.
\item The cell as ``first level of consciousness'' (Doc~170) is a working hypothesis, not a result of Doc~206.
\end{itemize}

\subsection{Relation within the corpus}

\begin{center}
\adjustbox{max width=\textwidth}{%
\begin{tabular}{c|p{8cm}|l}
\textbf{Document} & \textbf{Contribution to consciousness question} & \textbf{Status} \\
\hline
100 & Consciousness as fractal incoherence, Adlam connection & Connection paper \\
170 & Threshold ladder, $K_{\text{frak}}$ as threshold operator & Working hypothesis \\
193 & Non-closure theorem (anchored algebraically in 206) & Proven \\
206 & Z$_3$ skeleton, six bridges, $D_f = 3 - \xi$ & Proven \\
\textbf{207} & \textbf{Structural bridges between 206 and 100/170} & \textbf{Synthesis} \\
\end{tabular}%
}
\end{center}

Doc~207 makes no new statement about consciousness. It shows that the existing statements from Doc~100, Doc~170 and the IPI discussion sit within a consistent algebraic framework. The consciousness question itself remains an open research question. Doc~206 supplies the language in which further work can be precisely formulated; it does not supply an answer to what consciousness is.

\begin{philosophical}[Closing]
The question of what consciousness is is not answered by an algebraic bridge. But the question of in which language one can ask it without falling into numerological or merely metaphorical dead ends has, through Doc~206, gained sharper form. Nothing more is gained. Nothing less is intended.
\end{philosophical}

\section*{Acknowledgements}

The bridges formulated here build on the discussions in the \emph{Information Physics Institute} mailing list and on the work of Onur Tekermen, Phillips Paul, and Peter Austin. The phenomenological line owes itself to engagement with the no-go theorem of Adlam, McQueen and Waegell~\cite{Adlam2025}.

\begin{thebibliography}{9}
\bibitem{Adlam2025} E.\,C.~Adlam, K.\,J.~McQueen, M.~Waegell. \emph{Agency cannot be a purely quantum phenomenon.} arXiv:2510.13247 (2025).
\bibitem{Casali2013} A.\,G.~Casali, O.~Gosseries, M.~Rosanova, M.~Boly, S.~Sarasso, K.\,R.~Casali, S.~Casarotto, M.-A.~Bruno, S.~Laureys, G.~Tononi, M.~Massimini. \emph{A theoretically based index of consciousness independent of sensory processing and behavior.} \emph{Sci.\ Transl.\ Med.}\ \textbf{5}, 198ra105 (2013).
\bibitem{MaturanaVarela1980} H.\,R.~Maturana, F.\,J.~Varela. \emph{Autopoiesis and Cognition: The Realization of the Living.} D.~Reidel, Dordrecht (1980).
\bibitem{Margulis2001} L.~Margulis. \emph{The conscious cell.} \emph{Annals of the New York Academy of Sciences}\ \textbf{929}, 55--70 (2001).
\bibitem{Mancuso2018} S.~Mancuso, A.~Viola. \emph{Brilliant Green: The Surprising History and Science of Plant Intelligence.} Island Press (2018).
\bibitem{Trewavas2014} A.~Trewavas. \emph{Plant Behaviour and Intelligence.} Oxford University Press (2014).
\end{thebibliography}

\section*{License and availability}

\textbf{GitHub:} \href{https://github.com/jpascher/T0-Time-Mass-Duality}{\texttt{github.com/jpascher/T0-Time-Mass-Duality}}\\
\textbf{Zenodo (predecessor v1.0.13):} \href{https://zenodo.org/records/20041529}{\texttt{zenodo.org/records/20041529}} (DOI: \href{https://doi.org/10.5281/zenodo.20041529}{10.5281/zenodo.20041529})\\
\textbf{ORCID:} \href{https://orcid.org/0009-0000-6518-4064}{0009-0000-6518-4064}\\
\textbf{License:} CC-BY-SA 4.0
